% ============================================================
% PNAS RESEARCH ARTICLE — MAIN TEXT
% "Additive genetic channels: a dynamical theory of when
%  predictive regimes hold in evolution and breeding"
% Ortiz-Barrientos & Cooper
%
% TARGET: ~4,000 words body text, 40–50 references, 5 figures
% ARTICLE TYPE: EVOLUTION (or GENETICS)
% FORMAT: pnasresearcharticle (two-column, 9pt)
%
% COMPILATION
%   pdflatex pnas_research_article_main.tex
%   bibtex   pnas_research_article_main
%   pdflatex (x2)
% ============================================================

\documentclass[9pt,twocolumn,twoside]{pnas-new}

\articletype{EVOLUTION}

\templatetype{pnasresearcharticle}

% ----------------------------------------------------------
% GRAPHICS PATH
% All figure PDFs must be in the same directory as this .tex
% file, or the path below adjusted accordingly.
% PNAS production accepts PDF, EPS, and TIFF.
% ----------------------------------------------------------
\usepackage{graphicx}
\graphicspath{{./}}

% ----------------------------------------------------------
% CUSTOM MATH COMMANDS
% ----------------------------------------------------------
\newcommand{\vect}[1]{\boldsymbol{#1}}
\newcommand{\mat}[1]{\mathbf{#1}}
\newcommand{\Ag}{\mathcal{A}_{g}}
\newcommand{\Vlin}{V_{\mathrm{lin}}}
\newcommand{\Vquad}{V_{\mathrm{quad}}}
\newcommand{\CalL}{\mathcal{L}}
\newcommand{\tr}{\operatorname{tr}}
\newcommand{\Var}{\operatorname{Var}}
\newcommand{\Cov}{\operatorname{Cov}}
\newcommand{\Corr}{\operatorname{Corr}}

% ----------------------------------------------------------
% TITLE
% ----------------------------------------------------------
\title{Gene interactions are common yet additive prediction succeeds
       because evolution is local}

% ----------------------------------------------------------
% AUTHORS AND AFFILIATIONS
% ----------------------------------------------------------
\author[a,b,1]{Daniel Ortiz-Barrientos}
\author[b,c]{Mark Cooper}

\affil[a]{School of the Environment, The University of Queensland,
          St~Lucia, QLD~4072, Australia}
\affil[b]{ARC Centre for Plant Success in Nature and Agriculture,
          The University of Queensland, St~Lucia, QLD~4072, Australia}
\affil[c]{Queensland Alliance for Agriculture and Food Innovation,
          The University of Queensland, St~Lucia, QLD~4072, Australia}

\leadauthor{Ortiz-Barrientos}

% ----------------------------------------------------------
% SIGNIFICANCE STATEMENT  (50–120 words)
%
% AUDIENCE CALIBRATION NOTE:
%   This statement addresses any undergraduate-educated scientist.
%   It uses no technical vocabulary from quantitative genetics.
%   Its tone sets the register for the main text: biological paradox
%   first, mathematical resolution second, practical payoff third.
%   Every sentence in the Introduction should be writable by a reader
%   who found this statement convincing — not one who already knows
%   what a G matrix is.
% ----------------------------------------------------------
\significancestatement{%
Genes interact in complex networks, yet simple additive models predict evolutionary and agricultural change with surprising accuracy. We show why this works — and why it sometimes fails. A population does not experience an entire fitness landscape. It samples only the small region accessible through its current genetic variation, and in that region even a complex landscape can behave almost linearly. We derive a single index that measures how reliable additive prediction should be, and a rate diagnostic that detects whether reliability is improving or eroding before any drop in
forecast accuracy becomes visible. Applied to Darwin's finches and to breeding programs under climate change, the framework identifies which biological processes — directional selection, environmental rotation, recombination, and introgression — drive prediction failure and recovery.}

% ----------------------------------------------------------
% AUTHOR METADATA
% ----------------------------------------------------------
\authorcontributions{D.O.-B.\ and M.C.\ conceived the framework independently in their respective fields. D.O.-B.\ wrote the
first draft. M.C.\ contributed directly to framing and revision. Both authors contributed to interpretation and approved the final manuscript.}

\authordeclaration{The authors declare no competing interests.}

\correspondingauthor{\textsuperscript{1}Correspondence:
d.ortizbarrientos@uq.edu.au and mark.cooper@uq.edu.au}

% ----------------------------------------------------------
% KEYWORDS
% ----------------------------------------------------------
\keywords{fitness landscapes $|$ quantitative genetics $|$
          additive genetic variance $|$ adaptive evolution $|$
          genomic prediction $|$ channel switching}

% ----------------------------------------------------------
% ABSTRACT (≤ 250 words)
% ----------------------------------------------------------
\begin{abstract}
Additive genetic models predict evolutionary change and breeding response with surprising success despite pervasive gene interactions. We argue that this success is local. A population samples only the neighborhood of trait space defined by its additive genetic covariance matrix $\mat{G}$, and within that neighborhood, a curved fitness surface can often be approximated by its tangent plane. We define an additivity index $\Ag$ as the fraction of local heritable fitness variance captured by the gradient rather than by curvature --- a local $R^2$ for additive prediction. A logit-space rate diagnostic $r(t)$ then tracks whether a population is deepening into a predictive regime, approaching its boundary, or leaving it,  before the failure is visible in raw prediction accuracy. Applied to a one-trait moving-optimum model, the framework yields a counterintuitive result: a population near its fitness optimum is \emph{less} reliably additive than one far from it, because the directional signal collapses at the peak. Environmental rotation of the fitness gradient can exit an additive channel without any loss of genetic variance, a mechanism relevant to adaptation under climate change. In breeding, variance compression by directional selection deepens
a channel mechanically, and the direction of introgressed variance determines whether a wide cross helps or harms predictability. The framework identifies when additive prediction should hold, erode, fail, and recover, and which biological processes are responsible for these transitions, thereby converting a familiar first-order tool into a diagnostic of its own reliability.
\end{abstract}

\dates{This manuscript was compiled on \today}
\doi{\url{www.pnas.org/cgi/doi/10.1073/pnas.XXXXXXXXXX}}

\begin{document}

\maketitle
\thispagestyle{firststyle}
\ifthenelse{\boolean{shortarticle}}
  {\ifthenelse{\boolean{singlecolumn}}
    {\abscontentformatted}{\abscontent}}{}

\Firstpage

% ==========================================================
% INTRODUCTION
% (~800 words; no section heading in PNAS style)
%
% AUDIENCE NOTE: This Introduction is written for a PNAS reader
% who may be an ecologist, molecular biologist, or cell biologist
% encountering quantitative genetics from the outside. Technical
% terms are introduced with brief definitions on first use. The
% biological paradox is stated before the mathematical resolution.
% Classical tools are named and explained in one sentence each.
% ==========================================================

% ----------------------------------------------------------
% The paradox: paragraph 1
% ----------------------------------------------------------
Genes interact.
They act in pathways, share developmental machinery, and produce
fitness effects that depend on genetic background
\cite{Phillips2008,HuangMackay2016,SellaBarton2019}.
This complexity should make prediction hard.
Yet prediction is not hard, at least not always.
Breeders routinely use linear models to forecast the yield of
varieties never seen before \cite{Meuwissen2001,VanRaden2008}.
Evolutionary biologists describe the direction and speed of
adaptation using Lande's equation --- a model that treats all
genetic effects as additive --- even in populations where nonlinear
gene-by-gene interactions are well documented
\cite{Lande1979,WalshLynch2018}.
Something keeps the linear approximation alive in a deeply
nonlinear world.
The question is what.

% ----------------------------------------------------------
% Why this matters: paragraph 2
% ----------------------------------------------------------
The question is not merely academic.
When additive prediction breaks down, the consequences are
measurable.
A genomic breeding program that loses predictive accuracy wastes
selection cycles and misallocates resources
\cite{CooperMessina2023,Technow2021}.
A field study that applies Lande's equation outside its domain
of validity draws wrong conclusions about the direction or
magnitude of selection \cite{KingsolverPfennig2007}.
Climate change is rotating fitness landscapes --- shifting what
traits matter and how sharply fitness peaks --- at rates that
outpace the reorientation of many standing genetic variance
structures \cite{GrantGrant2002,GrantGrant2014}.
Knowing when additive prediction holds, and when it is about to
fail, is therefore a practical problem as much as a theoretical one.

\Endparasplit

% ----------------------------------------------------------
% What exists: classical tools, paragraph 3
% ----------------------------------------------------------
The theoretical toolkit for additive prediction rests on three
pillars.
Fisher's fundamental theorem identifies the additive component
of heritable fitness variance as the first-order driver of
mean fitness increase per generation \cite{Fisher1930,Price1972}.
Lande's multivariate response equation connects that additive
variance to the direction and magnitude of evolutionary change
\cite{Lande1979,LandeArnold1983}: the mean breeding value shifts
each generation in proportion to the product of the genetic
covariance matrix $\mat{G}$ --- which captures how heritable
variation is distributed and correlated across traits --- and the
fitness gradient $\vect{b}$, the direction in which selection
currently pushes.
Schluter's genetic lines of least resistance extend this logic to
show that the eigenvectors of $\mat{G}$ bias the trajectory of
adaptive radiation toward directions of abundant variance
\cite{Schluter1996,WalshBlows2009}.
Genomic prediction operationalizes the same idea at the level of
marker effects \cite{Meuwissen2001,VanRaden2008}.
All four tools share an assumption: that the local fitness surface,
as sampled by the current population, is well enough described by
its first-order (linear) term.

% ----------------------------------------------------------
% The gap: paragraph 4
% ----------------------------------------------------------
What is missing is a language for when that assumption holds,
when it is eroding, and when it recovers.
Nonlinear selection gradients --- the $\gamma$ matrix of
Lande and Arnold \cite{LandeArnold1983} and its extensions
\cite{PhillipsArnold1989} --- measure local curvature directly, but
they do not by themselves tell you how much that curvature matters
relative to the directional gradient across the population's
specific covariance cloud.
Fitness landscape theory \cite{WrightLandscape1932,Gavrilets2004}
describes global ruggedness but does not resolve the local
predictability question: a globally rugged landscape can support
a well-functioning additive channel if the population sits in a
locally flat valley.
A complete theory requires not just a static snapshot of local
geometry, but a longitudinal diagnostic that reveals whether
local additivity is strengthening or eroding over time.

% ----------------------------------------------------------
% The contribution: paragraph 5
% ----------------------------------------------------------
Here we provide both layers.
In earlier work, we defined the additivity index $\Ag$ as the
fraction of local heritable fitness variance explained by the
tangent-plane approximation --- a local $R^2$ for additive
prediction \cite{OrtizBarrientosCooper2026}.
High $\Ag$ identifies an \emph{additive channel}: a local state
in which first-order theory is informative.
The present paper extends this static picture to dynamics.
We derive a logit-space rate diagnostic $r(t)$ that decomposes the
gain or loss of local predictability into contributions from three
mechanistic routes: motion of the population mean across the
landscape, environmental forcing of local geometry, and
change in the covariance structure itself.
Within the covariance route, individual contributions from
selection, recombination, mutation, introgression, and drift
are explicitly identified.
The framework recovers Fisher's variance argument, Lande's response
equation, and genetic lines of least resistance as first-order
descriptions valid inside an additive channel.
It also predicts two qualitatively distinct failure modes ---
channel exit at the fitness optimum and channel exit from gradient
rotation --- and shows that these are separable from the sign
pattern of $r(t)$ before prediction accuracy visibly deteriorates.

% ----------------------------------------------------------
% Roadmap: paragraph 6
% ----------------------------------------------------------
We proceed as follows.
We first state the variance decomposition and define $\Ag$ (Results
section 1).
We then derive the dynamical closure, obtaining the transport laws
for gradient and curvature and the closed five-equation state-space
system (Results section 2).
We define $r(t)$ and its mechanistic decomposition (Results
section 3).
We apply these tools to two illustrative examples: the one-trait
moving-optimum model, which yields explicit formulas for $\Ag(t)$
and $r(t)$, and the two-trait gradient rotation model, which shows
how environmental change drives channel exit without any loss of
genetic variance (Results section 4).
We close by drawing out implications for natural selection,
adaptive radiation, and breeding under climate change.

% ==========================================================
\section*{Results}
% ==========================================================

% ----------------------------------------------------------
\subsection*{1.~The diagnostic layer: variance decomposition
             and the additivity index}
% ----------------------------------------------------------

% ----------------------------------------------------------
% The geometry: what b, H, and G mean
% ----------------------------------------------------------
Consider a population with mean breeding value $\bar{\vect{z}}$
--- the average additive genetic value across all traits ---
evolving under log-fitness surface $\ell(\vect{z})$.
Because the population's additive genetic covariance matrix
$\mat{G}$ defines the neighborhood of trait space that offspring
actually occupy, only the local geometry of $\ell$ matters for
short-term prediction.
Expand $\ell$ to second order around $\bar{\vect{z}}$:
\begin{equation}
\ell(\vect{z})
\approx
\ell(\bar{\vect{z}})
+ \vect{b}^\top\!\Delta\vect{z}
+ \tfrac{1}{2}\Delta\vect{z}^\top\!\mat{H}\,\Delta\vect{z},
\label{eq:taylor}
\end{equation}
where $\Delta\vect{z} = \vect{z} - \bar{\vect{z}}$, the fitness
gradient $\vect{b}$ points in the direction selection currently
pushes, and the Hessian matrix $\mat{H}$ describes how sharply
the fitness surface curves in each direction.
The Hessian is the local analogue of the Lande--Arnold $\gamma$
matrix of nonlinear selection gradients
\cite{LandeArnold1983,PhillipsArnold1989}.

% ----------------------------------------------------------
% Decomposition
% ----------------------------------------------------------
Under Gaussian breeding-value distributions, the local log-fitness
deviation $\delta\ell = L + Q$ (where $L = \vect{b}^\top\Delta\vect{z}$
is the linear part and $Q = \tfrac{1}{2}\Delta\vect{z}^\top
\mat{H}\,\Delta\vect{z}$ is the quadratic correction) decomposes
into variance components without a cross-term (SI Appendix,
Section~S1).
The cross-term vanishes by parity: $L$ changes sign under
$\Delta\vect{z}\mapsto-\Delta\vect{z}$ while $Q$ does not, so
symmetric points around the mean cancel in the product $LQ$.
The result is
\begin{equation}
\Var[\delta\ell] = \Vlin + \Vquad,
\label{eq:decomp}
\end{equation}
with gradient-aligned variance
$\Vlin = \vect{b}^\top\mat{G}\vect{b}$
and curvature-exposed variance
$\Vquad = \tfrac{1}{2}\tr[(\mat{H}\mat{G})^2]$.
$\Vlin$ involves one copy of $\mat{G}$; $\Vquad$ involves two.
That asymmetry --- one power versus two --- is the source of
the scaling results derived below.

% ----------------------------------------------------------
% Additivity index definition
% ----------------------------------------------------------
Their ratio defines the additivity index
\begin{equation}
\Ag = \frac{\Vlin}{\Vlin + \Vquad}.
\label{eq:Ag}
\end{equation}
$\Ag$ is the squared correlation between the linear predictor of
log-fitness and log-fitness itself: a local $R^2$ (SI Appendix,
Section~S1).
A population for which $\Ag > \alpha$ --- where $\alpha$ is a
threshold set by the precision required --- occupies an
\emph{additive channel} (Fig.~\ref{fig:regime}): a locally
predictive state in which Fisher's variance argument, Lande's
response equation, and genetic lines of least resistance are
reliable first-order descriptions, each nested inside the channel
condition (Fig.~\ref{fig:hierarchy}).
$\Ag$ depends jointly on the population's genetic structure,
the fitness gradient, and local curvature.
None of these three is sufficient alone.
A high-variance population can have low $\Ag$ if variance is spread
into directions of high curvature (Fig.~\ref{fig:concept}C);
a low-variance population in the same location may have high $\Ag$
because its narrow cloud sees a nearly flat slice of the landscape
(Fig.~\ref{fig:concept}B).

% ----------------------------------------------------------
% FIGURE 1 — inserted here, first cited above (Fig. 1B, 1C)
% 3-panel landscape composite, full two-column width.
% figure*[t!]: spans both columns; [t!] = top of next page only.
% NEVER use [b] or [bt] on figure* in twocolumn mode.
% ----------------------------------------------------------
\begin{figure*}[t!]
\centering
\includegraphics[width=17.8cm]{figure1_channel_concept}
\caption{\textbf{The additive channel concept.}
  The fitness landscape is a log-mixture-of-Gaussians surface
  $\ell(\mathbf{z}) = \log[\sum_k w_k \exp(-\frac{1}{2}
  (\mathbf{z}-\boldsymbol{\mu}_k)^\top
  \boldsymbol{\Sigma}_k^{-1}(\mathbf{z}-\boldsymbol{\mu}_k))]$;
  gradients and Hessians at each population position are computed
  by numerical differentiation, so all $\Ag$ annotations are exact.
  \textbf{(A)}~Two populations ($\mathcal{C}_1$, blue;
  $\mathcal{C}_2$, orange) traverse the landscape toward separate
  adaptive peaks.
  Ellipses show the additive genetic covariance cloud $\mat{G}$ at
  successive time points; $\Ag$ values annotate each step.
  As selection compresses $\mat{G}$, local predictability rises,
  regardless of global landscape ruggedness.
  \textbf{(B)}~A narrow cloud ($\sigma = 0.38$) at a fixed location.
  The green dashed tangent line tracks the true log-fitness surface
  (blue) closely across the sampled support; $\Ag \approx 0.80$.
  \textbf{(C)}~A broad cloud ($\sigma = 1.05$) at the same location.
  Orange shading between tangent and surface is the curvature load
  $\Vquad$ that the linear predictor misses; $\Ag \approx 0.34$.
  Expanding the cloud, not changing the landscape, halves the
  additivity index, illustrating the one-copy vs.\ two-copy
  scaling asymmetry ($\Vlin \propto G$; $\Vquad \propto G^2$).}
\label{fig:concept}
\end{figure*}

% ----------------------------------------------------------
% FIGURE 2 — inserted here, first cited above (Fig. 2)
% Single panel: Ag(s) rectangular hyperbola.
% Single-column figure[t]: does NOT span columns; places at top
% of next available column. Appropriate for single-panel figures
% whose content does not require full two-column width.
% Width = \columnwidth = 8.7 cm (PNAS single-column standard).
% ----------------------------------------------------------
\begin{figure}[t]
\centering
\includegraphics[width=\columnwidth]{figure2_regime_marker}
\caption{\textbf{$\Ag$ as a predictive regime marker.}
  Under uniform variance scaling $\mat{G}\to s\mat{G}$ with fixed
  $\vect{b}$ and $\mat{H}$, the additivity index follows
  $\Ag(s) = 1/(1 + s\kappa)$, where
  $\kappa = \Vquad_0/\Vlin_0$.
  Blue: interior channel ($\Ag > \alpha$).
  Amber: boundary zone ($|\Ag - \alpha| < 0.08$).
  Red: exterior ($\Ag < \alpha$).
  Wider clouds are always less locally additive, regardless of
  landscape shape.
  The state-space slice in the $(s,\kappa)$ plane is in
  SI Fig.~S2.}
\label{fig:regime}
\end{figure}

% ----------------------------------------------------------
% FIGURE 3 — inserted here, first cited above (Fig. 3)
% Square staircase diagram.
% Single-column figure[t]: nested-box diagram is square (5.2×5.4 in)
% and fits well at \columnwidth without spanning both columns.
% ----------------------------------------------------------
\begin{figure}[t]
\centering
\includegraphics[width=\columnwidth]{figure3_nested_hierarchy}
\caption{\textbf{Classical first-order results nest inside the
  additive channel.}
  Each concentric box is valid within the conditions at its corner.
  The channel condition $\Ag > \alpha$ (outermost, blue) is the
  prerequisite for all inner layers.
  Fisher's variance statement $\Vlin = \vect{b}^\top\mat{G}\vect{b}$
  (green) follows from Gaussian closure.
  Lande's response $\Delta\bar{\vect{z}} = \mat{G}\boldsymbol{\beta}$
  (orange) additionally requires $\mat{G}$ to be known and stable.
  Lines of least resistance $\vect{g}_{\max} = e_1(\mat{G})$
  (purple) further require $\mat{G}$ to persist across generations.
  The illustration shows misalignment angle $\theta$ between
  $\vect{b}$ (green) and $\vect{g}_{\max}$ (blue); the Lande
  response (orange dashed) lies between them.}
\label{fig:hierarchy}
\end{figure}

% ----------------------------------------------------------
\subsection*{2.~The dynamical closure: transport laws and the
             state-space system}
% ----------------------------------------------------------

% ----------------------------------------------------------
% Why dynamics are needed
% ----------------------------------------------------------
$\Ag$ is a snapshot.
Populations move, environments change, and genetic variance
structures reshape themselves under selection, recombination, and
the injection of new variation.
The channel that is open today may close tomorrow.
A warning that operates only at the moment of failure is no warning
at all.
What is needed is a diagnostic that reveals not just the current
regime, but the direction in which predictability is changing.

% ----------------------------------------------------------
% Logit transform and r(t)
% ----------------------------------------------------------
A logit transform $\CalL = \ln(\Vlin/\Vquad)$ converts the bounded
index $\Ag \in (0,1)$ to an unbounded scale suited for rate
analysis.
Its time derivative is the rate diagnostic
\begin{equation}
r(t) = \dot{\CalL}
     = \frac{\dot{\Vlin}}{\Vlin} - \frac{\dot{\Vquad}}{\Vquad},
\label{eq:r}
\end{equation}
which satisfies the exact algebraic identity
$\dot{\Ag} = \Ag(1-\Ag)\,r(t)$ (SI Appendix, Section~S2).
When $r > 0$, gradient-aligned variance is gaining on
curvature-exposed variance; when $r < 0$, it is losing ground.
A population can sit well inside a channel while $r < 0$: prediction
still works today, but is already being eroded.
That warning is precisely what a static snapshot cannot provide.

% ----------------------------------------------------------
% Transport laws
% ----------------------------------------------------------
Because $\vect{b}$ and $\mat{H}$ are partial derivatives of the
fitness surface evaluated at the moving population mean, their time
evolution follows from the chain rule alone (SI Appendix,
Section~S3):
\begin{align}
\dot{\vect{b}} &= \mat{H}\,\mat{G}\vect{b}
                + \vect{b}_u\,\dot{u},
\label{eq:bdot} \\
\dot{\mat{H}}  &= \mat{T}[\mat{G}\vect{b}]
                + \mat{H}_u\,\dot{u},
\label{eq:Hdot}
\end{align}
where $u$ is a scalar environmental state, $\vect{b}_u$ and
$\mat{H}_u$ are mixed derivatives capturing how the landscape
reshapes as the environment shifts, and $\mat{T}[\,\cdot\,]$
is a third-derivative tensor contraction.
The gradient changes for two reasons: the population mean moves
across a curved surface (the $\mat{H}\mat{G}\vect{b}$ term), or
the environment reshapes the surface at the current location
(the $\vect{b}_u\dot{u}$ term).
These two causes are biologically distinct and produce
distinguishably different patterns in $r(t)$.

% ----------------------------------------------------------
% Covariance dynamics and full state-space
% ----------------------------------------------------------
The companion covariance equation,
\begin{equation}
\dot{\mat{G}}
= \mat{G}\mat{H}\mat{G}
  - \lambda_{\mathrm{rec}}(\mat{G} - \mat{G}_*)
  + \mat{M}
  + \mat{J}(t)
  - \tfrac{1}{2N_e}\mat{G},
\label{eq:Gdot}
\end{equation}
shows that selection (the $\mat{G}\mat{H}\mat{G}$ term) deforms the
covariance cloud through the same local curvature that governs
mean displacement, while recombination ($\lambda_{\mathrm{rec}}$),
mutation ($\mat{M}$), introgression ($\mat{J}$), and drift
($1/2N_e$) each contribute independently
(SI Appendix, Section~S4).
Together, Eqs.~\eqref{eq:bdot}--\eqref{eq:Gdot} plus
$\dot{\bar{\vect{z}}} = \mat{G}\vect{b}$ form a closed five-equation
state-space system in which $\Ag$ and $r(t)$ are outputs, not inputs.

% ----------------------------------------------------------
\subsection*{3.~Mechanistic decomposition of $r(t)$}
% ----------------------------------------------------------

Differentiating $\Vlin$ and $\Vquad$ and substituting the transport
laws yields the three-way decomposition
$r = r_{\mathrm{motion}} + r_{\mathrm{env}} + r_{\mathrm{cov}}$,
where motion, environment, and covariance contributions are
separated explicitly (SI Appendix, Section~S5).
Within the covariance term, five biological processes each
contribute a signed component:
\begin{equation}
r_{\mathrm{cov}}
= r_{\mathrm{sel}}
+ r_{\mathrm{rec}}
+ r_{\mathrm{mut}}
+ r_{\mathrm{intro}}
+ r_{\mathrm{drift}}.
\label{eq:rcov}
\end{equation}
The drift term is always positive:
$r_{\mathrm{drift}} = +1/(2N_e)$.
This is a geometric consequence of the one-versus-two-copies
asymmetry in Eq.~\eqref{eq:decomp}: drift shrinks $\mat{G}$
uniformly, which reduces $\Vquad$ faster than $\Vlin$.
Small populations are therefore locally more additive than large
ones, not because their biology is simpler but because their narrow
covariance clouds see less curvature.
The mutation term is always negative: new variance enters
isotropically and increases curvature exposure faster than it
increases gradient alignment.
The selection, recombination, and introgression terms can each be
positive or negative, depending on whether the associated variance
change is better aligned with the gradient than with curvature.
The decomposition thus converts a scalar prediction failure into a
signed budget across five named biological processes.

% ----------------------------------------------------------
\subsection*{4.~Examples: three acts in the dynamics of
             local additivity}
% ----------------------------------------------------------

% ----------------------------------------------------------
% One-trait moving-optimum model
% ----------------------------------------------------------

% Act I and II
We illustrate the framework with two models of increasing
dimensionality.
Consider first a scalar trait under Gaussian stabilizing selection
toward a moving optimum $\theta(t)$:
\begin{equation}
\ell(z,t)
= \ell_0 - \tfrac{1}{2}s(z - \theta)^2,
\label{eq:landscape_1D}
\end{equation}
where $s > 0$ is the strength of stabilizing selection (the
curvature of the fitness peak) and $d = \theta - \bar{z}$ is the
lag between the optimum and the population mean.
At the current mean, the gradient $b = sd$ and the curvature
$H = -s$.
Substituting into Eq.~\eqref{eq:Ag}:
\begin{equation}
\Ag = \frac{d^2}{d^2 + G/2},
\label{eq:Ag_1D}
\end{equation}
where $G$ is the scalar genetic variance (SI Appendix, Section~S6).
This equation has a clear biological reading.
When the optimum is far away ($d^2 \gg G/2$), directional selection
dominates and $\Ag \approx 1$: the channel is wide open.
As the population closes the gap ($d \to 0$), the directional
signal $b = sd$ collapses, curvature dominates, and $\Ag \to 0$.
The channel closes not because the biology became complex but
because selection succeeded (Fig.~\ref{fig:dynamics}, Act~II).
A genomic breeder whose population approaches its target performance
ceiling will observe the same collapse in prediction accuracy for
the same reason.

The rate diagnostic in this model reduces to
\begin{equation}
r
= 2\frac{\dot{\theta}}{d}
  - sG
  + \lambda_{\mathrm{rec}}\!\left(1 - \frac{G_*}{G}\right)
  - \frac{m}{G}
  + \frac{1}{2N_e},
\label{eq:r_1D}
\end{equation}
where $m$ is the per-generation mutational variance input and
$G_*$ is the recombination baseline (SI Appendix, Section~S6).
Each term in Eq.~\eqref{eq:r_1D} is directly interpretable.
Faster optimum motion relative to current lag ($2\dot{\theta}/d$)
regenerates directional signal.
Wider variance ($-sG$) increases curvature exposure.
Mutation ($-m/G$) and introgression add variance isotropically
and always erode local additivity.
Drift ($+1/2N_e$) increases local additivity even while reducing
evolvability.
The three phases of Fig.~\ref{fig:dynamics} --- optimum moves,
stops, moves again --- illustrate channel entry, exit, and
re-entry as the sign of $r(t)$ oscillates between positive and
negative.

% ----------------------------------------------------------
% FIGURE 4 (moving optimum) — inserted here, first cited above.
% 3-panel landscape 1×3 composite, full two-column width.
% figure*[t!]: must be [t!] — never [b] on figure* in twocolumn.
% ----------------------------------------------------------
\begin{figure*}[t!]
\centering
\includegraphics[width=17.8cm]{figure5_moving_optimum}
\caption{\textbf{Channel dynamics in the one-trait
  moving-optimum model.}
  Parameters: $s = 0.10$, $\lambda_{\mathrm{rec}} = 0.30$,
  $G_* = 0.80$, $m = 0.005$, $N_e = 500$.
  Green background: optimum moving ($\dot\theta = 0.02$).
  Orange background: optimum stationary ($\dot\theta = 0$).
  \textbf{(A)}~Lag $d(t)$ (green) and variance $G(t)$ (blue dashed)
  on the left axis; additivity index $\Ag(t)$ (purple) on the right.
  Channel opens as lag grows (Act~I), closes as population nears
  optimum (Act~II), reopens when optimum moves again (Act~I).
  \textbf{(B)}~Logit trajectory $\CalL(t)$; entry, exit, and re-entry
  annotated at threshold $\CalL_\alpha$.
  Slope of $\CalL(t)$ equals $r(t)$ exactly.
  \textbf{(C)}~Rate diagnostic $r(t)$ (black) decomposed into
  $r_{\mathrm{motion}} = 2\dot\theta/d$ (green) and
  $r_{\mathrm{curv}} = -sG$ (orange).
  Gap marks $|d| < \delta$.
  Positive $r(t)$ precedes rising $\Ag$; negative $r(t)$ precedes
  falling $\Ag$ — the diagnostic warns before the failure is visible.}
\label{fig:dynamics}
\end{figure*}

% Act III: gradient rotation
The moving-optimum model holds the gradient direction fixed and
varies only its magnitude.
A qualitatively different failure mode arises when the gradient
rotates without any change in variance magnitude or curvature
magnitude --- the scenario most relevant to climate-driven
environmental change.

Let the fitness gradient $\vect{b}$ rotate by angle $\theta$ away
from the major axis $\vect{g}_{\max} = e_1(\mat{G})$ of the
covariance ellipsoid (Fig.~\ref{fig:rotation}A).
In two traits with eigenvalues $\lambda_1 > \lambda_2$ of
$\mat{G}$:
\begin{equation}
\Vlin(\theta)
= \|\vect{b}\|^2(\lambda_1\cos^2\!\theta + \lambda_2\sin^2\!\theta),
\label{eq:Vlin_rot}
\end{equation}
while $\Vquad = \tfrac{1}{2}\tr[(\mat{H}\mat{G})^2]$ is
independent of gradient direction because it depends on the
eigenvalues of $\mat{H}$ and $\mat{G}$ but not their relative
orientation (Fig.~\ref{fig:rotation}B).
$\Ag$ therefore declines monotonically as the gradient rotates
toward the minor axis of $\mat{G}$: channel exit from misalignment,
not from reaching the optimum.

This is \textit{Act~III} (Fig.~\ref{fig:rotation}C), and it is
separable from Act~II by the rate diagnostic.
In Act~II, $r_{\mathrm{env}} \approx 0$ and $r_{\mathrm{cov}} < 0$:
the channel closes because selection has compressed the covariance
cloud toward the now-irrelevant directional axis.
In Act~III, $r_{\mathrm{env}} \ll 0$ while $r_{\mathrm{cov}}$ can
remain positive if selection is still compressing the cloud along
the \emph{old} gradient direction.
That sign pattern --- strongly negative environmental contribution
combined with positive or zero covariance contribution --- does not
arise in Act~II.
The two scenarios may be invisible in raw prediction accuracy at the
moment of channel exit, but are distinguishable from $r(t)$ before
that moment.

% ==========================================================
\section*{Discussion}
% ==========================================================

% ----------------------------------------------------------
% What Ag measures that heritability does not
% ----------------------------------------------------------
The additivity index $\Ag$ does not replace heritability.
It answers a different question.
Heritability asks how much phenotypic variance is genetic.
$\Ag$ asks, of the heritable fitness variance, how much does the
linear approximation capture.
A highly heritable trait in a population sitting on a fitness peak
can have $\Ag \approx 0$: the variance is real, but none of it is
captured by the gradient because the gradient has collapsed.
A moderately heritable trait in a population far from its optimum
can have $\Ag \approx 1$: the tangent-plane approximation is
excellent, and Lande's equation will predict response faithfully.
The two indices together --- heritability and $\Ag$ --- give a
fuller picture than either alone.

% ----------------------------------------------------------
% Channel switching in natural populations
% ----------------------------------------------------------
The three-act framework connects classical results in evolutionary
biology to a common mechanistic language.
Darwin's finches on Daphne Major illustrate all three acts within
a decade \cite{GrantGrant2002,GrantGrant2014}.
The 1977 drought imposed a steep directional gradient on beak
depth --- Act~I: channel open, beak depth shifted nearly a
millimetre in one generation.
As the population tracked the new optimum across the 1980s, the
lag shrank and the channel narrowed --- Act~II.
El Niño years subsequently rotated the seed community toward small
soft seeds, rotating the fitness gradient away from the major axis
of beak variance --- Act~III.
The rate diagnostic $r(t)$ would have distinguished Acts~II and~III
in real time: the sign of $r_{\mathrm{env}}$ is negative in Act~III
but near zero in Act~II.
More broadly, the framework suggests that adaptive radiations into
new ecological opportunity \cite{Schluter2000} are episodes of
channel opening: a population encounters a new directional gradient
that is well-aligned with its standing genetic variance, $\Ag$
rises, and first-order evolution proceeds rapidly along genetic
lines of least resistance \cite{Schluter1996}.
Conversely, evolutionary stasis may reflect not absence of genetic
variance but absence of a directional gradient: $\Ag$ is low
because the gradient has collapsed, not because $\mat{G}$ is small.

% ----------------------------------------------------------
% FIGURE 5 (gradient rotation) — inserted here, cited above.
% 3-panel landscape 1×3 composite, full two-column width.
% ----------------------------------------------------------
\begin{figure*}[t!]
\centering
\includegraphics[width=17.8cm]{figure4_gradient_rotation}
\caption{\textbf{Gradient rotation drives channel exit without
  loss of genetic variance (Act~III).}
  Parameters: $\mat{G} = \mathrm{diag}(\lambda_1,\lambda_2)$,
  $\mat{H} = -\gamma\mat{I}$ ($\gamma = 0.414$), $|\vect{b}| = 1$.
  \textbf{(A)}~Fixed covariance ellipse under three gradient
  directions $\vect{b}_1$ (0°), $\vect{b}_2$ (50°),
  $\vect{b}_3$ (90°) relative to $\vect{g}_{\max} = e_1(\mat{G})$.
  $\Ag$ values are exact; $\Vquad$ is constant across all $\theta$.
  \textbf{(B)}~$\Ag(\theta)$ for four eccentricities
  $\lambda_1/\lambda_2$.
  $\Vquad$ constant; $\Vlin$ falls with rotation.
  Eccentricity $\lambda_1/\lambda_2 = 2$ is the marginal case,
  just touching $\alpha = 0.70$ at $\theta = 90°$.
  \textbf{(C)}~Exit angle $\theta^*$ vs curvature-to-gradient
  ratio $\kappa_0$.
  All curves converge at $\kappa_0 = (1-\alpha)/\alpha \approx 0.43$.
  Environmental change that rotates the fitness gradient erodes the
  channel silently: no variance loss, no curvature increase.}
\label{fig:rotation}
\end{figure*}
% ----------------------------------------------------------
% Breeding implications
% ----------------------------------------------------------
For breeding programs, the framework converts three qualitative
intuitions into quantitative predictions.
First, directional selection compresses $\mat{G}$, which deepens
a channel mechanically: $\Vquad \propto G^2$ falls faster than
$\Vlin \propto G$, so a more compressed cloud is locally more
additive.
A breeding program that has successfully selected toward a target
environment will therefore find genomic models increasingly reliable
--- until the population nears its local optimum and Act~II begins.
Second, the introgression criterion derived in the SI Appendix
separates benefit from cost in a wide cross:
\[
\delta\Ag \propto
\frac{(\vect{b}^\top\vect{v})^2}{\Vlin}
-
\frac{\vect{v}^\top\mat{H}\mat{G}\mat{H}\vect{v}}{\Vquad},
\]
where $\vect{v}$ is the direction of incoming variance.
A wide cross helps predictability when the first term dominates
--- when the new variance is gradient-aligned --- and harms it
when the second term dominates.
This criterion can guide germplasm choice in programs facing
yield plateaus \cite{Technow2021}.
Third, and most urgently under climate change: the rate diagnostic
tells you \emph{why} $\Ag$ is falling.
Act~II and Act~III require opposite interventions.
Act~II calls for environmental redesign --- finding new selective
targets, expanding the breeding population's effective lag.
Act~III calls for genetic redesign --- introducing variance
aligned with the new gradient, through either recombination or
targeted introgression.
Measuring $\vect{b}$, $\mat{H}$, and $\mat{G}$ is demanding in
natural populations \cite{Kingsolver2001,KingsolverPfennig2007}
but tractable in long-term multi-environment trials and experimental
evolution systems, where the framework converts a familiar
first-order tool into a diagnostic for its own reliability.

% ----------------------------------------------------------
% Limits and open problems
% ----------------------------------------------------------
Three limits of the current framework deserve emphasis.
First, the variance decomposition rests on Gaussian closure.
Non-Gaussian breeding-value distributions --- expected under strong
selection, low heritability, or major-effect loci --- will produce
a nonzero cross-term $\Cov(L,Q)$, and $\Ag$ will no longer equal
the local $R^2$ exactly.
The framework gives qualitative guidance in those cases but
loses its $R^2$ interpretation.
Second, $\mat{H}$ involves third-order derivatives of the fitness
surface when the population mean is moving (Eq.~\ref{eq:Hdot}):
the dynamical theory of curvature is one derivative order above
the dynamical theory of the gradient.
In practice, estimating third derivatives requires large, carefully
designed experiments.
Third, the state-space system is open with respect to the fitness
surface itself: $\ell(\vect{z},u)$ must be supplied from outside.
The framework is therefore diagnostic and predictive conditional
on a known or estimated landscape, not a self-contained
evolutionary model.
Closing the loop --- deriving the landscape from an explicit
ecological model and carrying it through time alongside the
population --- is the natural next step and the subject of
continuing work.

% ==========================================================
% DATA AVAILABILITY
% ==========================================================
\dataavail{Code used to generate all figures is available at
[repository DOI to be inserted before submission].
No new empirical data were analyzed in this theoretical study.}

% ==========================================================
% ACKNOWLEDGMENTS
% ==========================================================
\acknow{[Funding and acknowledgment information to be added before
submission.]}

\showacknow{}

% ==========================================================
% BIBLIOGRAPHY
% \bibsplit[n] splits the reference list across two columns
% at reference n. With ~40 references, split at 20.
% ==========================================================

\bibsplit[20]
\bibliography{ortiz-barrientos_cooper_research_article_references}

\end{document}